\documentclass[11pt]{article}

\usepackage[a4paper,margin=25mm]{geometry}
\usepackage{amsmath,amssymb,mathtools}
\usepackage{booktabs,longtable,array}
\usepackage{enumitem}
\usepackage{xcolor}
\usepackage{hyperref}
\usepackage{listings}
\usepackage{microtype}

\hypersetup{
  colorlinks=true,
  linkcolor=blue!55!black,
  urlcolor=blue!65!black,
  pdftitle={Error, Norm, and Convergence-Rate Computation in the Distributed Periodic Maxwell Study},
  pdfauthor={DiscoGMPI}
}

\newcommand{\vect}[1]{\boldsymbol{#1}}
\newcommand{\norm}[1]{\left\lVert #1 \right\rVert}
\newcommand{\abs}[1]{\left| #1 \right|}
\newcommand{\code}[1]{\nolinkurl{#1}}
\newcommand{\mcode}[1]{\mathtt{\detokenize{#1}}}
\newcommand{\R}{\mathbb{R}}
\newcommand{\Oh}{\Omega_h}
\newcommand{\Eex}{\vect E^{\mathrm{ex}}}
\newcommand{\Hex}{\vect H^{\mathrm{ex}}}
\newcommand{\Eh}{\vect E_h}
\newcommand{\Hh}{\vect H_h}

\lstdefinestyle{shell}{
  basicstyle=\ttfamily\small,
  frame=single,
  breaklines=true,
  columns=fullflexible,
  showstringspaces=false
}

\title{Error, Norm, and Convergence-Rate Computation\\
in the Distributed Periodic Poisson-Bracket Maxwell Study}
\author{DiscoGMPI Technical Report}
\date{June 15, 2026}

\begin{document}

\maketitle
\tableofcontents
\newpage

\section{Purpose and scope}

This report documents, in mathematical and implementation terms, the
diagnostics computed by the convergence driver
\begin{center}
\code{examples/convergence_distributed_periodic_poisson_bracket_maxwell.jl}.
\end{center}
It covers:
\begin{itemize}[leftmargin=*]
  \item component-wise errors for
        \(E_x,E_y,E_z,H_x,H_y,H_z\);
  \item aggregate electric, magnetic, and six-component field errors;
  \item relative field errors;
  \item quadrature-sampled \(L^\infty\) diagnostics;
  \item energy discrepancy and discrete energy drift;
  \item mesh-size measures;
  \item observed convergence rates;
  \item the final PASS/FAIL convergence verdict.
\end{itemize}

The report follows the current source code. It also identifies distinctions
that are important when interpreting the CSV output:
\begin{enumerate}[leftmargin=*]
  \item The CSV distinguishes the requested and measured mesh sizes.
        \code{target_h} records the Gmsh request. Rates use the measured
        \code{characteristic_h}.
  \item \code{energy_error} compares numerical and analytical energies at the
        final time, while \code{relative_energy_error} measures numerical
        energy drift from the initial numerical energy.
  \item the reported \(L^\infty\) values are maxima over cubature points, not
        mathematically exact suprema over the continuous domain.
  \item component rates for exact-zero components measure decay of numerical
        leakage and can become dominated by roundoff.
\end{enumerate}

\section{Relevant source-code path}

The computation is split across the following files.

\begin{longtable}{>{\raggedright\arraybackslash}p{0.38\textwidth}
                  >{\raggedright\arraybackslash}p{0.56\textwidth}}
\toprule
Source & Responsibility \\
\midrule
\endhead
\code{examples/convergence_distributed_periodic_poisson_bracket_maxwell.jl}
&
Mesh-family generation, time integration, component errors, sampled
\(L^\infty\) errors, rate calculation, CSV writing, and PASS/FAIL logic.
\\
\code{examples/distributed_periodic_poisson_bracket_maxwell.jl}
&
Analytical plane wave, nodal-to-cubature interpolation, aggregate field
norms, analytical energy, charge, and related quadrature diagnostics.
\\
\code{examples/meshes/periodic_box_unstructured.geo}
&
Periodic box geometry and the Gmsh parameter
\code{NxTarget}, with \(h_{\mathrm{target}}=L_x/N_x\).
\\
\code{src/solver/kernels/metrics.jl}
&
Affine reference-to-physical tetrahedron map and element Jacobian.
\\
\code{src/solver/kernels/reference_tet.jl}
&
Degree-\(N\) nodal tetrahedral basis, Vandermonde matrices, mass matrix, and
reference derivative operators.
\\
\code{src/solver/kernels/JaskowiecSukumar.jl}
&
Jaskowiec--Sukumar tetrahedral cubature points and weights.
\\
\code{src/DistributedDG.jl}
&
Distributed mass-matrix Maxwell energy and MPI reduction over owned
elements.
\\
\bottomrule
\end{longtable}

\section{Analytical Maxwell problem}

\subsection{Material model and wave parameters}

The convergence driver currently enforces
\[
  \varepsilon=1,\qquad \mu=1,
\]
although the analytical helper functions retain the general symbols
\(\varepsilon\) and \(\mu\). The wave speed and dispersion relation are
\[
  c=\frac{1}{\sqrt{\varepsilon\mu}},
  \qquad
  \omega = kc = \frac{k}{\sqrt{\varepsilon\mu}}.
\]
The magnetic amplitude is \(H_0=1\).

Let \(x_0=x_{\min}\) be the lower \(x\)-coordinate inferred from the mesh.
The phase is
\[
  \theta(x,t)=k(x-x_0)-\omega t.
\]
The analytical fields used for initialization and final-time error
measurement are
\begin{align}
  \Eex(x,y,z,t)
    &=
    \left(
      0,\,
      0,\,
      -\frac{kH_0}{\omega\varepsilon}\sin\theta
    \right),
  \\
  \Hex(x,y,z,t)
    &=
    \left(
      0,\,
      H_0\sin\theta,\,
      0
    \right).
\end{align}
For the default \(\varepsilon=\mu=H_0=1\), \(\omega=k\), so
\[
  \Eex=(0,0,-\sin\theta),
  \qquad
  \Hex=(0,\sin\theta,0).
\]

The exact nonzero components are therefore \(E_z\) and \(H_y\). The exact
components \(E_x,E_y,H_x,H_z\) are identically zero. This fact is important
when interpreting their component-wise errors and rates.

\subsection{Periodic compatibility}

For a box of length \(L_x=x_{\max}-x_{\min}\), the driver requires
\[
  \frac{kL_x}{2\pi}\in\mathbb{N}.
\]
Thus an integer number of wavelengths fits across the periodic \(x\)
direction. Opposite faces in all three coordinate directions are paired,
using translations
\[
  (-L_x,0,0),\qquad
  (0,-L_y,0),\qquad
  (0,0,-L_z).
\]

\section{Mesh family and mesh-size definitions}

\subsection{Gmsh target refinement}

The Gmsh geometry defines
\[
  h_{\mathrm{target}}=\frac{L_x}{N_x},
\]
where \(N_x\) is the value of \code{NxTarget}. With a base target
\(N_{x,0}\), the convergence driver constructs
\[
  N_{x,\ell}=2^\ell N_{x,0},
  \qquad \ell=0,1,\ldots,L-1.
\]
Therefore,
\[
  h_{\mathrm{target},\ell}
    =\frac{L_x}{N_{x,\ell}}
    =2^{-\ell}h_{\mathrm{target},0}.
\]
For the default four levels, this is exactly
\[
  h_0,\quad \frac{h_0}{2},\quad \frac{h_0}{4},\quad \frac{h_0}{8}.
\]

Gmsh creates a new unstructured tetrahedral mesh for every level. The meshes
are not required to be nested. \code{target_h} represents the requested
Gmsh size, not the diameter of every tetrahedron.

\subsection{Characteristic mesh size}

The driver also reports
\[
  h_{\mathrm{char}}
    =
    \left(\frac{\abs{\Omega}}{N_{\mathrm{el}}}\right)^{1/3},
\]
where \(N_{\mathrm{el}}\) is the global tetrahedron count. Element volume is
calculated from the affine map:
\[
  \abs{K}
    =
    \abs{\widehat K}\abs{\det J_K},
  \qquad
  \abs{\widehat K}=\frac{4}{3}.
\]
Each MPI rank sums only its owned element volumes. An
\code{MPI.Allreduce(...,+,...)} produces the global volume and global element
count.

\paragraph{Important distinction.}
The convergence-rate denominator uses the measured, volume-derived size stored
in \code{characteristic_h}. The requested Gmsh size is retained separately as
\code{target_h}. On an unstructured Gmsh family the characteristic size need
not decrease by exactly a factor of two, even though the requested target size
does.

\section{DG representation and evaluation at cubature points}

\subsection{Nodal polynomial field}

On each tetrahedron \(K\), every scalar component is represented by a
degree-\(N\) nodal polynomial:
\[
  u_h|_K(\vect x)
    =
    \sum_{j=1}^{N_p}u_{K,j}\,
    \ell_j\!\left(F_K^{-1}(\vect x)\right),
\]
with
\[
  N_p=\frac{(N+1)(N+2)(N+3)}{6}.
\]
The six component arrays are
\[
  E_x,\ E_y,\ E_z,\ H_x,\ H_y,\ H_z
  \in \R^{N_p\times N_{\mathrm{local}}}.
\]

\subsection{Reference tetrahedron and physical map}

The reference tetrahedron has vertices
\[
  (-1,-1,-1),\quad
  (1,-1,-1),\quad
  (-1,1,-1),\quad
  (-1,-1,1).
\]
For reference coordinates \((r,s,t)\), the affine shape functions are
\begin{align}
  \phi_1 &= -\frac{r+s+t+1}{2},&
  \phi_2 &=  \frac{r+1}{2},\\
  \phi_3 &=  \frac{s+1}{2},&
  \phi_4 &=  \frac{t+1}{2}.
\end{align}
The physical point is
\[
  \vect x=F_K(r,s,t)
    =\sum_{a=1}^{4}\phi_a(r,s,t)\vect x_a.
\]
The Jacobian is constant on each affine tetrahedron:
\[
  J_K=
  \begin{bmatrix}
    (\vect x_2-\vect x_1)/2 &
    (\vect x_3-\vect x_1)/2 &
    (\vect x_4-\vect x_1)/2
  \end{bmatrix}.
\]
Consequently,
\[
  d\vect x=\abs{\det J_K}\,d\vect \xi.
\]

\subsection{Interpolation matrix}

The solution is stored at DG interpolation nodes, while errors are evaluated
at volume cubature points. The code constructs
\[
  I_q=V_qV^{-1},
\]
where \(V\) is the nodal Vandermonde matrix and \(V_q\) evaluates the same
orthonormal modal basis at cubature points. If
\(\vect u_K=(u_{K,1},\ldots,u_{K,N_p})^T\), then
\[
  u_h(\vect x_{K,q})=(I_q)_{q,:}\vect u_K.
\]
This interpolation is performed independently for all six field components.

\section{Volume cubature}

For DG order \(N\), the convergence driver selects the
Jaskowiec--Sukumar cubature order
\[
  p_q=\max(2,\,2N+4).
\]
Let \(\{(\vect \xi_q,w_q)\}_{q=1}^{N_q}\) be the resulting rule on
\(\widehat K\). A physical-domain integral is approximated by
\[
  \int_K f(\vect x)\,d\vect x
  \approx
  \sum_{q=1}^{N_q}
    \abs{\det J_K}\,w_q\,f(F_K(\vect \xi_q)).
\]
The implementation calls
\[
  W_{K,q}=\abs{\det J_K}\,w_q
\]
the \emph{physical weight}.

The cubature rule exactly integrates polynomials up to its supported degree.
However, the analytical wave contains sine functions. Therefore,
quantities such as
\((u_h-u^{\mathrm{ex}})^2\) are not polynomials, and the reported norm is a
high-order quadrature approximation to the continuous norm rather than a
symbolically exact integral. The choice \(2N+4\) deliberately exceeds the
degree needed for the square of a degree-\(N\) polynomial and reduces
quadrature contamination from the trigonometric analytical field.

\section{\texorpdfstring{Component-wise \(L^2\) errors}{Component-wise L2 errors}}

For each scalar component \(u\in
\{E_x,E_y,E_z,H_x,H_y,H_z\}\), define
\[
  e_u(\vect x,t_f)=u_h(\vect x,t_f)-u^{\mathrm{ex}}(\vect x,t_f).
\]
The continuous component norm is
\[
  \norm{e_u}_{L^2(\Omega)}
  =
  \left(
    \sum_{K\in\mathcal{T}_h}
    \int_K e_u^2\,d\vect x
  \right)^{1/2}.
\]
The implemented cubature approximation is
\[
  \norm{e_u}_{L^2(\Omega),q}
  =
  \left[
    \sum_{K\in\mathcal{T}_h}
    \sum_{q=1}^{N_q}
    W_{K,q}
    \left(
      (I_q)_{q,:}\vect u_K
      -
      u^{\mathrm{ex}}(\vect x_{K,q},t_f)
    \right)^2
  \right]^{1/2}.
\]

The corresponding CSV fields are:
\begin{center}
\begin{tabular}{lll}
  \code{l2_ex_error},&
  \code{l2_ey_error},&
  \code{l2_ez_error},\\
  \code{l2_hx_error},&
  \code{l2_hy_error},&
  \code{l2_hz_error}.
\end{tabular}
\end{center}

\subsection{Exact-zero components}

For this plane wave,
\[
  E_x^{\mathrm{ex}}=E_y^{\mathrm{ex}}
  =H_x^{\mathrm{ex}}=H_z^{\mathrm{ex}}=0.
\]
Thus, for example,
\[
  \norm{e_{E_x}}_{L^2}
  =\norm{E_{x,h}}_{L^2}.
\]
These norms measure spurious cross-component leakage. On an unstructured
tetrahedral mesh, interpolating a function of \(x\) into the full
three-dimensional degree-\(N\) polynomial space does not generally produce a
polynomial that depends on \(x\) alone. Its transverse derivatives can
therefore generate nonzero inactive field components during evolution. Those
components can converge one order more slowly than the active components and
can eventually dominate an aggregate vector norm. Their observed rates can
also become unstable once their magnitudes approach floating-point roundoff.
The physically active component errors
\(\norm{e_{E_z}}_{L^2}\) and \(\norm{e_{H_y}}_{L^2}\) are generally the most
direct component-level measures of plane-wave accuracy.

\section{Aggregate electric, magnetic, and total errors}

At a point, the vector errors are
\[
  \vect e_E=\Eh-\Eex,\qquad
  \vect e_H=\Hh-\Hex.
\]
Their squared Euclidean magnitudes are
\begin{align}
  \abs{\vect e_E}^2
    &=
    e_{E_x}^2+e_{E_y}^2+e_{E_z}^2,\\
  \abs{\vect e_H}^2
    &=
    e_{H_x}^2+e_{H_y}^2+e_{H_z}^2.
\end{align}
The aggregate errors are
\begin{align}
  \mathcal{E}_E
    &=\norm{\Eh-\Eex}_{L^2(\Omega)},\\
  \mathcal{E}_H
    &=\norm{\Hh-\Hex}_{L^2(\Omega)},\\
  \mathcal{E}_{EH}
    &=
    \left(
      \mathcal{E}_E^2+\mathcal{E}_H^2
    \right)^{1/2}.
\end{align}
Because all terms use the same quadrature,
\begin{align}
  \mathcal{E}_E^2
    &=
    \norm{e_{E_x}}_{L^2,q}^2+
    \norm{e_{E_y}}_{L^2,q}^2+
    \norm{e_{E_z}}_{L^2,q}^2,\\
  \mathcal{E}_H^2
    &=
    \norm{e_{H_x}}_{L^2,q}^2+
    \norm{e_{H_y}}_{L^2,q}^2+
    \norm{e_{H_z}}_{L^2,q}^2.
\end{align}

The CSV names are:
\[
  \mcode{l2_electric_error}=\mathcal{E}_E,\qquad
  \mcode{l2_magnetic_error}=\mathcal{E}_H,\qquad
  \mcode{l2_total_error}=\mathcal{E}_{EH}.
\]

\section{Relative total field error}

The analytical combined field norm is
\[
  \mathcal{N}_{EH}^{\mathrm{ex}}
  =
  \left(
    \norm{\Eex}_{L^2,q}^2+
    \norm{\Hex}_{L^2,q}^2
  \right)^{1/2}.
\]
The CSV quantity \code{relative_total_error} is
\[
  \mathcal{E}_{EH,\mathrm{rel}}
  =
  \frac{\mathcal{E}_{EH}}
       {\max(\mathcal{N}_{EH}^{\mathrm{ex}},
             \epsilon_{\mathrm{mach}})}.
\]
The machine-epsilon guard prevents division by zero. For the nontrivial plane
wave the denominator is not close to zero.

\section{\texorpdfstring{Quadrature-sampled \(L^\infty\) errors}
       {Quadrature-sampled Linfinity errors}}

The code also computes
\begin{align}
  \mathcal{E}_{E,\infty}^{(q)}
    &=
    \max_{K,q}
    \abs{\vect e_E(\vect x_{K,q},t_f)},\\
  \mathcal{E}_{H,\infty}^{(q)}
    &=
    \max_{K,q}
    \abs{\vect e_H(\vect x_{K,q},t_f)},\\
  \mathcal{E}_{EH,\infty}^{(q)}
    &=
    \max_{K,q}
    \sqrt{
      \abs{\vect e_E(\vect x_{K,q},t_f)}^2+
      \abs{\vect e_H(\vect x_{K,q},t_f)}^2
    }.
\end{align}
Each rank computes local maxima over owned elements. An MPI maximum reduction
produces the global result.

The CSV fields are
\[
  \mcode{linf_electric_error},\quad
  \mcode{linf_magnetic_error},\quad
  \mcode{linf_total_error}.
\]

\paragraph{Interpretation.}
These are not exact continuous \(L^\infty\) norms because extrema between
cubature points are not searched. A more precise name would be
``cubature-point maximum error.'' The values become more representative as
the mesh and sampling rule are refined, but they should not be presented as
rigorous suprema.

\section{Energy diagnostics}

\subsection{Analytical and numerical energy at final time}

The electromagnetic energy density is
\[
  w(\vect x,t)
  =
  \frac{1}{2}\varepsilon\abs{\vect E}^2+
  \frac{1}{2}\mu\abs{\vect H}^2.
\]
The quadrature diagnostic computes
\begin{align}
  \mathcal{H}_h(t_f)
    &\approx
    \sum_{K,q}W_{K,q}
    \left[
      \frac{1}{2}\varepsilon\abs{\Eh(\vect x_{K,q},t_f)}^2+
      \frac{1}{2}\mu\abs{\Hh(\vect x_{K,q},t_f)}^2
    \right],\\
  \mathcal{H}^{\mathrm{ex}}(t_f)
    &\approx
    \sum_{K,q}W_{K,q}
    \left[
      \frac{1}{2}\varepsilon\abs{\Eex(\vect x_{K,q},t_f)}^2+
      \frac{1}{2}\mu\abs{\Hex(\vect x_{K,q},t_f)}^2
    \right].
\end{align}
The CSV field \code{energy_error} is the signed difference
\[
  \mcode{energy_error}
  =
  \mathcal{H}_h(t_f)-\mathcal{H}^{\mathrm{ex}}(t_f).
\]
It is not an absolute value and is not normalized.

\subsection{Discrete energy drift}

The convergence driver separately evaluates the numerical energy with the DG
mass matrix:
\[
  \mathcal{H}_{M,h}
  =
  \frac{1}{2}
  \sum_{K}
  \abs{\det J_K}
  \left[
    \varepsilon
    \sum_{\alpha=x,y,z}
      \vect E_{\alpha,K}^{\,T}M\vect E_{\alpha,K}
    +
    \mu
    \sum_{\alpha=x,y,z}
      \vect H_{\alpha,K}^{\,T}M\vect H_{\alpha,K}
  \right].
\]
The CSV field \code{relative_energy_error} is
\[
  \mcode{relative_energy_error}
  =
  \frac{
    \mathcal{H}_{M,h}(t_f)-\mathcal{H}_{M,h}(0)
  }{
    \max(\mathcal{H}_{M,h}(0),\epsilon_{\mathrm{mach}})
  }.
\]
It measures conservation of the numerical Hamiltonian. It is not the
relative version of \code{energy_error}. The current CSV naming should be
interpreted with this distinction in mind.

\section{Charge diagnostics}

The source-free exact solution has zero electric and magnetic charge. The
code evaluates the DG polynomial divergences and integrates
\begin{align}
  Q_E
    &=
    \int_\Omega
      \varepsilon\,\nabla\cdot\Eh\,d\vect x,\\
  Q_H
    &=
    \int_\Omega
      \mu\,\nabla\cdot\Hh\,d\vect x.
\end{align}
The quadrature approximations are written as
\code{electric_charge} and \code{magnetic_charge}. These are global signed
integrals, not \(L^2\) norms of divergence. Positive and negative local
divergence errors can cancel.

\section{MPI ownership and reductions}

Every volume loop is restricted to
\code{distributed_mesh.partition.owned}. Ghost elements are excluded from
integration. This prevents interface halo copies from being counted more
than once.

The reduction pattern is:
\begin{itemize}[leftmargin=*]
  \item squared-error integrals, energies, volume, element counts, and charges:
        \code{MPI.Allreduce(local,+,comm)};
  \item sampled \(L^\infty\) values:
        \code{MPI.Allreduce(local,max,comm)};
  \item stable time-step estimate:
        \code{MPI.Allreduce(local_dt,min,comm)};
  \item measured elapsed time:
        \code{MPI.Allreduce(local_elapsed,max,comm)}.
\end{itemize}

All ranks therefore obtain the same global diagnostics. The CSV is written
only by rank zero.

\section{Time integration and its effect on measured error}

For polynomial order \(N\), the convergence test uses an H-first explicit
symplectic partitioned Runge--Kutta scheme of order
\[
  p_t=N+1.
\]
Each rank estimates a stable local step. The global estimate is the minimum
over ranks. The final number of steps and actual step are
\[
  n_{\mathrm{step}}
    =
    \max\left(
      1,\,
      \left\lceil
        \frac{T}{\Delta t_{\mathrm{est}}}
      \right\rceil
    \right),
  \qquad
  \Delta t=\frac{T}{n_{\mathrm{step}}}.
\]
This guarantees that the numerical solution lands exactly at the requested
final time \(T\).

The measured field error is the combined result of:
\begin{enumerate}[leftmargin=*]
  \item spatial DG discretization error;
  \item time-integration error;
  \item volume-cubature error in the reported norm;
  \item floating-point and MPI reduction error.
\end{enumerate}
The choice \(p_t=N+1\) and a mesh-dependent CFL step is intended to prevent a
low-order temporal method from dominating the spatial study. Nevertheless,
a strict spatial-convergence investigation should repeat selected cases with
a smaller CFL and verify that rates and errors are insensitive to further
time-step reduction.

\section{Observed convergence rates}

\subsection{Rate formula}

For two consecutive mesh levels, let \(e_c\) be the coarse error and \(e_f\)
the fine error. The code computes
\[
  p_{\mathrm{obs}}
  =
  \frac{
    \log(e_c/e_f)
  }{
    \log(h_c/h_f)
  }.
\]
Here
\[
  h_c=\mcode{characteristic_h}_{c},
  \qquad
  h_f=\mcode{characteristic_h}_{f}.
\]
The measured ratio \(h_c/h_f\) is read from the two result rows. It is not
assumed to equal two, because an independently generated unstructured Gmsh
family generally follows the requested target sizes only approximately.

The first mesh level has no preceding level, so its rate is stored as missing
and written as an empty CSV field. If either error is nonpositive, the helper
also returns missing. In normal finite-precision calculations the norms are
nonnegative and are positive unless the error is represented as exactly zero.

\subsection{Rates stored in the CSV}

The aggregate rates are:
\[
  \mcode{rate_electric},\qquad
  \mcode{rate_magnetic},\qquad
  \mcode{rate_total}.
\]
The component rates are:
\begin{center}
\begin{tabular}{lll}
  \code{rate_ex},&
  \code{rate_ey},&
  \code{rate_ez},\\
  \code{rate_hx},&
  \code{rate_hy},&
  \code{rate_hz}.
\end{tabular}
\end{center}
Every rate is computed independently from the corresponding error column.

\subsection{Meaning of rates for zero exact components}

Rates for \(E_x,E_y,H_x,H_z\) describe how rapidly spurious components decay.
While the leakage is resolved, these rates may be closer to \(N\) than to
\(N+1\) and can lower the aggregate field rate. If the errors reach roundoff,
their ratios can instead become erratic, extremely large, negative, or
missing. The console prints \(E_z\), \(H_y\), and transverse-leakage
diagnostics so that these regimes can be distinguished.

\section{PASS/FAIL criterion}

For each DG order \(N\), the verdict uses only the final rate from the finest
two mesh levels:
\[
  p_E^{\mathrm{final}}
  =\mcode{rate_electric}\text{ on the finest row},
\]
\[
  p_H^{\mathrm{final}}
  =\mcode{rate_magnetic}\text{ on the finest row}.
\]
The strict target rates encoded in the driver are
\[
  p_E^{\mathrm{expected}}=N+1,
  \qquad
  p_H^{\mathrm{expected}}=N.
\]
With the configured tolerance \(0.5\), an order passes if
\begin{align}
  p_E^{\mathrm{final}}&\ge N+1-0.5,\\
  p_H^{\mathrm{final}}&\ge N-0.5.
\end{align}
Both values must be present and finite. The global verdict is PASS only if
both tests pass for every requested DG order.

The \(E=N+1\) check is an enhanced aggregate-field target inherited from the
aligned Poisson-bracket convergence experiment. It is not a generic theorem
for every independently generated unstructured tetrahedral family. In
particular, order-\(N\) transverse polarization leakage can make the full
electric vector converge more slowly even when the active \(E_z\) component
converges at order \(N+1\). A failed strict verdict must therefore be
interpreted together with the component table rather than as evidence that
the error integration is incorrect.

The following quantities do \emph{not} currently affect PASS/FAIL:
\begin{itemize}[leftmargin=*]
  \item \code{rate_total};
  \item all six component rates;
  \item monotonicity at every intermediate level;
  \item relative field error;
  \item \(L^\infty\) diagnostics;
  \item energy discrepancy or energy drift;
  \item charge diagnostics.
\end{itemize}

\section{CSV data dictionary}

\small
\begin{longtable}{>{\raggedright\arraybackslash}p{0.29\textwidth}
                  >{\raggedright\arraybackslash}p{0.65\textwidth}}
\toprule
Column & Definition \\
\midrule
\endhead
\code{mpi_ranks} & Number of MPI ranks in the run. \\
\code{boundary_condition} & Always \code{periodic} for this driver. \\
\code{order} & DG polynomial degree \(N\). \\
\code{esprk_order} & Time-integrator order \(N+1\). \\
\code{cubature_order} & Jaskowiec--Sukumar order \(\max(2,2N+4)\). \\
\code{mesh_level} & Zero-based refinement level \(\ell\). \\
\code{nx_target} & Gmsh \code{NxTarget} value \(N_{x,\ell}\). \\
\code{target_h} & \(L_x/N_{x,\ell}\); requested Gmsh target size. \\
\code{nelements} & Global tetrahedron count. \\
\code{min_owned_elements} & Minimum owned-element count over ranks. \\
\code{max_owned_elements} & Maximum owned-element count over ranks. \\
\code{characteristic_h} &
\((\abs{\Omega}/N_{\mathrm{el}})^{1/3}\); measured size used for rates.
\\
\code{dt} & Actual step \(T/n_{\mathrm{step}}\). \\
\code{nsteps} & Number of time steps. \\
\code{elapsed_seconds} & Maximum time-integration wall time over ranks. \\
\code{l2_electric_error} & \(\norm{\Eh-\Eex}_{L^2,q}\). \\
\code{l2_magnetic_error} & \(\norm{\Hh-\Hex}_{L^2,q}\). \\
\code{l2_total_error} &
\(\sqrt{\mathcal{E}_E^2+\mathcal{E}_H^2}\).
\\
\code{l2_ex_error} & \(\norm{E_{x,h}-E_x^{\mathrm{ex}}}_{L^2,q}\). \\
\code{l2_ey_error} & \(\norm{E_{y,h}-E_y^{\mathrm{ex}}}_{L^2,q}\). \\
\code{l2_ez_error} & \(\norm{E_{z,h}-E_z^{\mathrm{ex}}}_{L^2,q}\). \\
\code{l2_hx_error} & \(\norm{H_{x,h}-H_x^{\mathrm{ex}}}_{L^2,q}\). \\
\code{l2_hy_error} & \(\norm{H_{y,h}-H_y^{\mathrm{ex}}}_{L^2,q}\). \\
\code{l2_hz_error} & \(\norm{H_{z,h}-H_z^{\mathrm{ex}}}_{L^2,q}\). \\
\code{relative_total_error} &
\(\mathcal E_{EH}/\mathcal N_{EH}^{\mathrm{ex}}\), with an epsilon guard.
\\
\code{linf_electric_error} & Maximum electric vector error at cubature points. \\
\code{linf_magnetic_error} & Maximum magnetic vector error at cubature points. \\
\code{linf_total_error} & Maximum combined six-component error at cubature points. \\
\code{energy_error} &
Signed final quadrature energy difference
\(\mathcal H_h(t_f)-\mathcal H^{\mathrm{ex}}(t_f)\).
\\
\code{relative_energy_error} &
Relative mass-matrix energy drift from \(t=0\) to \(t_f\).
\\
\code{electric_charge} &
\(\int_\Omega\varepsilon\nabla\cdot\Eh\,d\vect x\), by cubature.
\\
\code{magnetic_charge} &
\(\int_\Omega\mu\nabla\cdot\Hh\,d\vect x\), by cubature.
\\
\code{rate_electric} & Consecutive-level rate of \code{l2_electric_error}. \\
\code{rate_magnetic} & Consecutive-level rate of \code{l2_magnetic_error}. \\
\code{rate_total} & Consecutive-level rate of \code{l2_total_error}. \\
\code{rate_ex}, \code{rate_ey}, \code{rate_ez} &
Consecutive-level rates of the three electric component errors.
\\
\code{rate_hx}, \code{rate_hy}, \code{rate_hz} &
Consecutive-level rates of the three magnetic component errors.
\\
\bottomrule
\end{longtable}
\normalsize

\section{Consistency identities useful for auditing results}

The following identities should hold up to floating-point summation error:
\begin{align}
  \mcode{l2_electric_error}^2
  &\approx
  \mcode{l2_ex_error}^2+
  \mcode{l2_ey_error}^2+
  \mcode{l2_ez_error}^2,\\
  \mcode{l2_magnetic_error}^2
  &\approx
  \mcode{l2_hx_error}^2+
  \mcode{l2_hy_error}^2+
  \mcode{l2_hz_error}^2,\\
  \mcode{l2_total_error}^2
  &\approx
  \mcode{l2_electric_error}^2+
  \mcode{l2_magnetic_error}^2.
\end{align}
These are particularly useful for detecting component ordering mistakes or
inconsistent analytical fields.

For the specified wave, one also expects:
\begin{itemize}[leftmargin=*]
  \item \(E_z\) and \(H_y\) dominate the field and error norms;
  \item \(E_x,E_y,H_x,H_z\) remain small;
  \item electric and magnetic aggregate errors are often similar because the
        exact wave has equal electric and magnetic energy for
        \(\varepsilon=\mu=1\);
  \item electric and magnetic charge integrals remain close to zero;
  \item relative energy drift remains small for the symplectic method.
\end{itemize}

\section{Recommended execution and interpretation workflow}

\subsection{Run the convergence study}

\begin{lstlisting}[style=shell]
mpiexec -n 2 julia --project=. \
  examples/convergence_distributed_periodic_poisson_bracket_maxwell.jl \
  --nx-target=2 \
  --levels=4 \
  --orders=2,3,4 \
  --periods=1 \
  --cfl=0.05 \
  --output=output/convergence_distributed_periodic_poisson_bracket.csv
\end{lstlisting}

\subsection{Interpretation order}

For each DG order:
\begin{enumerate}[leftmargin=*]
  \item verify that \code{target_h} halves and that
        \code{characteristic_h} decreases consistently at every level;
  \item verify that aggregate and active-component errors decrease;
  \item check the three norm-consistency identities above;
  \item inspect \(E_z\) and \(H_y\) rates;
  \item inspect aggregate electric and magnetic rates;
  \item treat zero-component rates as leakage diagnostics;
  \item verify that a smaller CFL does not materially change the finest
        errors or rates;
  \item inspect energy drift and charge as conservation checks;
  \item use the automated PASS/FAIL verdict only after these checks.
\end{enumerate}

\section{Limitations and interpretation cautions}

\begin{enumerate}[leftmargin=*]
  \item \textbf{Non-nested unstructured meshes.}
        The Gmsh family follows a uniform target-size sequence but is not a
        nested finite-element hierarchy. Pre-asymptotic rate variation is
        therefore expected.
  \item \textbf{Target size versus actual sizes.}
        The requested value \(L_x/N_x\) is a mesher-control parameter. Rates
        use the measured volume-based \code{characteristic_h}, avoiding the
        false assumption that Gmsh exactly halves the realized element size.
  \item \textbf{Space--time error coupling.}
        The final error includes temporal error. A CFL sensitivity check is
        necessary before attributing all observed behavior to spatial order.
  \item \textbf{Quadrature approximation.}
        Trigonometric analytical fields make the norm integrands
        non-polynomial. The selected high-order cubature is accurate but not
        algebraically exact.
  \item \textbf{Sampled \(L^\infty\).}
        The maximum is taken only at volume cubature points.
  \item \textbf{Charge cancellation.}
        Integrated signed charge can be small even when local divergence
        errors are not.
  \item \textbf{Zero-component rates.}
        Exact-zero components can contain order-\(N\) polarization leakage
        before eventually becoming roundoff-dominated. They contribute to the
        aggregate vector norms.
  \item \textbf{Energy column naming.}
        \code{energy_error} and \code{relative_energy_error} use different
        energy comparisons, as documented above.
  \item \textbf{Verdict coverage.}
        PASS/FAIL checks only the finest aggregate electric and magnetic
        rates against a strict enhanced-order target. It does not certify all
        columns or all mesh transitions, and the electric target is not
        guaranteed on a generic unstructured tetrahedral family.
\end{enumerate}

\section{Summary}

The convergence driver evaluates the degree-\(N\) DG solution at
Jaskowiec--Sukumar volume cubature points, maps those points to each physical
tetrahedron, compares all six numerical components with the analytical plane
wave, and integrates squared errors using
\(\abs{\det J_K}w_q\). MPI sums owned-element contributions globally.

Component norms are combined by Euclidean sums to produce electric, magnetic,
and total field norms. Consecutive-level rates use the measured
characteristic-size ratio and are written for aggregate and component errors.
The automated verdict uses only the finest aggregate electric and magnetic
rates, requiring at least \(N+1-0.5\) and \(N-0.5\), respectively. The console
also reports the active \(E_z\) and \(H_y\) errors and transverse leakage,
because the strict aggregate electric target can be reduced by order-\(N\)
polarization leakage on unstructured tetrahedra.

The methodology is suitable for distributed \(h\)-convergence assessment
provided that users also check temporal sensitivity, quadrature adequacy,
pre-asymptotic mesh behavior, and the special interpretation required for
exact-zero field components.

\end{document}
